
\newcommand\w{5}
\section{Result}

\begin{Form}
\noindent Commit ID: \vspace{3mm} 

\TextField[
  name=commitid,
  multiline=true,
  width=10cm,
  height=1cm
]{}

\end{Form}\\

\bigskip

\noindent Which calibration points were used. Position 1-10 belongs to zone 1, 11 - 20 to zone 2 etc. Fill in a checkbox even if not all positions in the zone was used.
\bigskip

\CheckBox[name=Zone 1,charsize=12pt]{Zone 1:}\quad
\CheckBox[name=Zone 2,charsize=12pt]{Zone 2:}\quad
\CheckBox[name=Zone 3,charsize=12pt]{Zone 3:}\quad
\CheckBox[name=Zone 4,charsize=12pt]{Zone 4:}

\bigskip






\begin{Form}
\begin{table}[h!]
\renewcommand{\arraystretch}{1.6}
\centering

\begin{tabular}{|l|p{4cm}|p{7cm}|}
\hline
\textbf{Metric} & \textbf{Value} & \textbf{Description} \\
\hline

RMSE &
\TextField[name=RMSE, width=4cm]{} &
Root Mean Square Error, the difference between the robots coordinates and the cameras coordinates after transformation to the robots frame.(mm) \\
\hline

STD &
\TextField[name=STD, width=4cm]{} &
Standard deviation error between robot coordinates and transformed camera coordinates.(mm) \\
\hline

Number of calibration positions &
\TextField[name=CalibrationPositions, width=4cm]{} &
How many calibration positions were used. \\
\hline

Camera distance to robot base &
\TextField[name=CameraDistance, width=4cm]{} &
Measured distance between camera, x-y plane, and robot base, absolute value.(mm)\\
\hline

Camera height &
\TextField[name=CameraHeight, width=4cm]{} &
Measured height between camera and the robot base, absolute value.(mm)\\
\hline

Estimated camera position &
\TextField[name=EstimatedCameraPosition, width=4cm]{} &
The rough angle of the cameras position to the robot base, e.g. 150 degrees, left of robot. (degrees)\\
\hline

Estimated camera pose &
\TextField[name=EstimatedCameraPose, width=4cm]{} &
The estimated angle of camera to the robot base, e.g. z = -20, y = -10 (degrees)\\
\hline



\end{tabular}

\caption{Result table}
\end{table}
\end{Form}












\begin{comment}
    


\begin{Form} % Start interactive form environment
\begin{table}[h!]

\renewcommand{\arraystretch}{3} % Increase row height
\centering

\begin{tabular}{|l|l|l|}
\hline

\textbf{Metric} & \textbf{Description} & \textbf{Value} \\
\hline

Accuracy & \TextField[name=accuracy_disc,
  multiline=true, width=\w cm]{} & 
  \TextField[name=accuracy,
  multiline=true, width=\w cm]{} \\
Percentage correct & \TextField[name=percentage_disc,
  multiline=true, width=\w cm]{} & \TextField[name=percentage,
  multiline=true, width=\w cm]{} \\
Amount of tries & \TextField[name=tries_disc,
  multiline=true, width=\w cm]{} & \TextField[name=tries,
  multiline=true, width=\w cm]{} \\
Total score & \TextField[name=score_disc,
  multiline=true, width=\w cm]{} & \TextField[name=score,
  multiline=true, width=\w cm]{} \\
\hline
\end{tabular}
\caption{Result table}
\end{table}
\end{Form}

\end{comment}